\chapter{54}
  
\section{Dienstag, 28. August} 

Das Fernglas lag im Regal neben dem Küchenschrank. Kurz vor sechs
stand Louis beim gedeckten Frühstückstisch vor der Hütte. Mit
ruhiger Bewegung liess er das Gerät den Weg entlang gleiten, der den
gegenüberliegenden Berghang durchquerte. Weit entfernt schliesslich entdeckte er einen
kleinen Punkt in Bewegung. Der vergrösserte sich. Mazi würde bald da
sein.
Louis hörte den Kaffee hochplubbern und eilte in die Hütte. Er verehrte
die einfachen Espressokocher. Der hier oben musste bereits ein langes
Leben hinter sich haben. Das Aluminium der kleinen Maschine hatte sich längst
angebräunt. Der Drehverschluss sperrte. Womöglich lag es an der
Dichtung. Worauf es aber wirklich ankam, war der Duft nach würzigem
Kaffee. Wohlbekannt. Louis schloss die Augen. Es hatte sich von selbst
ergeben. Nach seinem Zusammenzug mit Giorgia hatte er seit Beginn die
Rolle des Kaffeebereiters. Das Gebräu hatte sie morgens am
wirkungsvollsten aus den Federn geholt. Er seufzte. Die Erinnerung
stimmte ihn traurig.
Mazis Stimme stellte Louis zurück in die Küchenecke vor den Herd.

«Hi, Ciro, das sieht unglaublich aus.»

Mit einem Hopser landete Mazi in der Hütte. Er trat neben Louis und klopfte
ihm auf die Schulter.

«Danke, hi, du kommst gerade richtig, der Kaffee ist fertig.»

Louis wies Mazi nach draussen. Sie setzten sich gegenüber an den Tisch.
So flink wie sich Mazi bewegte, so besonnen sprach er. Dabei hing sein
Blick ungebrochen an Louis’ Augen.
Mazi war ein Erzähler. Er sprach ruhig und ohne Punkt und Komma. Auf
eine Sprechpause würde Louis lange warten. Der Typ atmete ganz
offensichtlich während des Redens. Louis schöpfte sich gerade eine
weitere Portion Rührei, als Mazi seine Gabel zum ersten Mal belud.
Unterdessen kannte Louis den Tagesplan und wusste, dass sie bald
losgehen würden.

«Okay für dich, Ciro?»

Mazi schaute Louis kurz an und begann zu essen. 

«Schmeckt vorzüglich.»

\sectionbreak

Sie durchquerten die hüglige Landschaft, als Mazi seine Hände als
Trichter vor seinen Mund legte.

«Flexa, back!»

Neben einem grossen Gesteinsbrocken schoss nun ein kleines Tier
pfeilschnell hoch. Sie blieben stehen.
Mazi ging in die Knie. Louis wunderte sich, wie schnell die Hündin auf sie zu gerannt kam
und sich alsbald hechelnd neben Mazi stellte. Der begann sie zu
streicheln und sanft abzuklopfen.

«Braver Hund.»

Flexa wedelte und gab heisere Laute von sich. Die Szene wirkte wie ein
frohes Wiedersehen nach einer langen Trennung. Louis beobachtete die zwei mit Vergnügen.

«Take time!»

Nach Mazis Befehl blieb Flexa augenblicklich stehen. Unbeweglich schaute
sie zu ihm hoch.
Während sie das letzte Stück bis zum Weideplatz gingen, nutzte Mazi die
Zeit. Er legte Louis seine Verbindung zu Flexa dar.

«Sie kam vor drei Jahren zu Joh. Ich war damals das erste Mal als
Mitarbeiter hier oben auf der Alp. Joh bezog mich in die Dressur mit
ein. So ergab es sich, dass sie lernte, uns beide als Meister zu
akzeptieren. Selbst wenn wir zusammen anwesend sind, versteht Flexa, wer
gerade das Sagen hat. Sie ist eine tolle Hirtenhündin. Ein kluges Tier.»

Als sie ihren Namen hörte, blickte die Hündin treuherzig zu Mazi auf.

«Du wirst ohne sie auskommen. Solltest du aber ein Schaf vermissen,
rufst du uns und ich komme mit Flexa zu dir rüber. Sie holt jeden
Ausreisser zurück.»

Mazi lachte und streichelte der Hündin über den Kopf.
Nach wie vor hielt Louis Ausschau nach den Schafen. Weit und breit kein
Tier. Er schaute ratlos um sich.
Als könne er Gedanken lesen, blieb Mazi stehen. Er streckte seinen Arm
aus und zeigte in Richtung der verschiedenen, kleinen Hügel und
Felsbrocken an der Steilwand vor ihnen.

«Dort sind sie, Ciro,» er fuhr mit dem Finger der Wand entlang, «unsere
Schafe haben sich bereits wieder versteckt. Sie mögen die intensive
Bergsonne nicht. Sie sind vorwiegend in den Dämmerungsstunden aktiv.
Oder ganz früh am Morgen. Das Schwarznasenschaf ist ein ruhiges Tier.
Das karge Futterangebot reicht ihm, es ist standorttreu, geländegängig,
sehr robust und eben, genügsam,» nun legte Mazi seine Hand auf Louis’
Schulter, «damit unterscheidet sich die Schafrasse von andern, zum
Beispiel vom weissen Alpenschaf. Das kennst du ja.»

Louis nickte. Mazis Erklärungen klangen schlüssig. Nur wusste Louis noch
immer nicht, was er ausser dem Einfangen ausgebüxter Tiere den ganzen
Tag hier oben zu tun hatte.

«Joh und Manuela sömmern ihre Schafe in Geländekammern und in Gruppen,
wie du weisst. Die Tiere kennen ihre Alp bestens, müssen hier oben nicht
eingezäunt werden und eigentlich auch nicht jeden Moment behirtet, weil
sie, wie gesagt, standorttreu sind.»

«Mhm,» Louis hüstelte, «verstehe ich. Nur, worin genau liegt denn unsere
Aufgabe mit den Tieren?»

Jetzt fiel Mazi wieder in seinen atemlosen Vortragsmodus.

«Ich verstehe deine Frage. Also, deine Aufgaben bestehen darin, dass du
die Schafe jeden Morgen und Abend zählst. Jedes Tier spürst du einmal
täglich auf und schaust es dir genau an. Ist es gesund und unverletzt?
Wenn ihm etwas fehlt, rufst du mich an. Sollte ich nicht gerade
erreichbar sein, meldest du dich bei Manuela. Sie wird dich beraten, was
zu tun ist oder gegebenenfalls hochkommen. Sie verarztet die Tiere,
sofern nötig. Manchmal gibt es auch Probleme mit trächtigen Schafen. Die
werfen im Herbst ihre Lämmer. Diesbezüglich rufst du Manuela lieber
einmal mehr an als zu spät.
Und ja, es kommt vor, dass es Wolfsalarm gibt, das hatten wir diesen
Sommer erst ein Mal. Wir haben die Aufgabe, die Tiere zu beschützen.
Flexa ist dann besonders auf Trab. Manchmal fordern wir zusätzliche
Hunde an. Im schlimmsten Fall unterstützen uns die Jäger. Geschossen
werden darf der Wolf nicht, erschreckt aber schon,» Mazi vergrösserte
die Augen, «und ein spezieller Grund für die Behirtung der Schafe ist
Johs Integrationsprogramm für die Jungs. Davon hat dir Joh erzählt,
oder?»

Mazi kramte eine Tüte mit Nüssen aus seiner Hosentasche.

«Ja,» Louis fuhr sich über die Stirn, «sie sollen hier oben zur Ruhe
kommen, eine Art Kulturwechsel erleben, weg von ihrem alten Umfeld, oder?»

Louis stockte. Eigentlich war er eins zu eins einer von denen. Mit
dem Unterschied, inoffiziell hier zu sein. Und dabei musste es bleiben. 
Für einen Moment wünschte er, seine Geschichte würde
auf dem Tisch liegen. Ende der Aufführung. Ende des Versteckspiels. 
Er schüttelte unmerklich den Kopf. Was für ein Gedanke.
Mazi breitete seine Arme aus, die Handflächen nach oben.

«Exakt, weg von all den Versuchungen. Und die Hirtenaufgaben sind ein
wichtiger Teil von Johs Konzept.»

Mazis Unterkunft lag wenige hundert Meter weiter unten. Nach wie vor im Schatten
des Morgens. Auch der heutige Tag versprach, heiss zu werden.
Die Hütte war klein. Sie bestand aus einem einzigen Raum. Mit einem
Blick ins Innere stillte sich Louis’ Neugierde. Die Unordnung fiel ihm als Erstes auf. 
Benutztes Geschirr stand auf dem kleinen Holztisch am Fenster, Jacken und Schuhe lagen da und dort,
und neben einem breiten, zerwühlten Bett schien ein Haufen schmutziger
Wäsche zu lagern.

«Schau nicht genau hin, Ciro,» Mazi lachte laut, «mein Haushalt liegt im
Argen. Ich musste gestern dringend eine Studie zu Ende schreiben und
rausschicken.»

Bevor sie sich Louis’ Schafen in der Nähe seiner Alphütte zuwenden
wollten, setzten sie sich auf eine verwitterte Bank inmitten eines Flecks dürrer Wiese.
Louis hatte seine Neugierde bisher zügeln können. Mazis Herkunft
interessierte ihn. Und was hatte es wohl mit diesem afrikanischen Song
auf sich, der Joh verzaubert hatte?

«Woher kommst du ursprünglich?»

Was Louis hier oben und in seiner persönlichen Situation zweifellos
lernen konnte, war zuzuhören.

«Meine Eltern begegneten sich an der Uni. Sie studierten
\emph{Internationale Beziehungen} in Genf. Papa ist Afrikaner. Mama eine
Bernerin. Die verliebten sich und da gab es mich draus. Als Jugendliche
waren Mama und Joh befreundet, also so eine lose Sache, wie Mama sagt,»
Mazi zog die Schultern hoch, hob die Arme und grinste, «die haben den
Kontakt nie ganz abgebrochen. Auch Papa und Joh verstehen sich prima.
Ich war ein paar Mal bei Joh zu Besuch, ferienhalber. Joh rief uns dann
Jahre später um Hilfe an, weil er wusste, dass ich als Agro-Ingenieur
arbeite. Er hatte gerade alle Hände voll zu tun mit einem seiner
Schützlinge. Ich sollte Manuela entlasten. So bin ich nun das dritte
Jahr für einige Wochen hier, um auszuhelfen und meine Studien
voranzubringen. Ich arbeite an einem Entwicklungsprojekt in Südafrika.»

Mazi schraubte den Deckel auf die Trinkflasche.

«Interessant, und an dieser Studie hast du gestern gearbeitet?»

Louis machte sich auf einen weiteren Vortrag gefasst. Die Vorstellung
hatte etwas Aufregendes. Mazis Worte klangen auffallend
selbstverständlich. Je länger er ihm zuhörte, desto wohler fühlte er
sich. Der Kerl strahlte eine beeindruckende Zufriedenheit aus. Der stand mit beiden Beinen am
Boden. Selbst seine Bewegungen sahen mühelos aus. Mazi streckte sich.

«Hm, genau, und heute Abend erzählst du von dir, okay?»

Er hob den Daumen.

Louis musste sich zwingen, ihn anzusehen.

«Ja, klar.»

Er war erleichtert, dass sich Mazi bereits abgewandt hatte. Immerhin,
Louis hatte den ganzen Tag Zeit, sich eine passende Geschichte
auszudenken.
